% GENERATED FILE. Edit canonical lessons, then run npm run generate:books.
% canonical-source-hash: fnv1a64:d96ea2b185fcd9cd
% canonical-lessons: MR-C155-jodne, MR-C155-bandhne, MR-C155-sivne, MR-C155-rangavne, MR-C155-kadhne

\chapter{To join, To tie, To sew, To paint, To take out}
\label{ch:mr-to-join-to-tie-to-sew-to-paint-to-take-out}
\hlchaptermodality{155}

\begin{chapteropening}
\textbf{By the end of this chapter:} \emph{I can say to join, to tie, to sew, to paint and to take out.}

Complete the last lesson of chapter 155: i can say to join, to tie, to sew, to paint and to take out.
\end{chapteropening}

\section[joḍṇe]{\mr{जोडणे} (joḍṇe) — to join}

\label{lesson:MR-C155-jodne}



\begin{quote}
Before the new one: say the Marathi for to break off, then the Marathi for to break.
\end{quote}

\subsection*{You'll want to know: \mr{जोडणे}}
\textbf{\mr{जोडणे}} — \emph{joḍṇe} — \textquotedblleft{}to join\textquotedblright{}.

\begin{grammarlens}[title={What you've built}]
One more word for this chapter.
\end{grammarlens}

\subsection*{Your turn}
\begin{itemize}
\raggedright
  \item \emph{Say it:} \emph{joḍṇe}
  \item \emph{Say it:} \emph{joḍṇe}, once more
  \item \emph{Say it:} say \emph{joḍṇe}
  \item \emph{Recall:} say the Marathi for to break off, then the Marathi for to break, then say \emph{joḍṇe} again
\end{itemize}

\begin{tcolorbox}[breakable,colback=teal!4,colframe=teal!35!black,title={Before you move on}]
What does \mr{जोडणे} mean? (\textquotedblleft{}To join\textquotedblright{}.) Say it once more.
\end{tcolorbox}
\section[bāndhṇe]{\mr{बांधणे} (bāndhṇe) — to tie, to build}

\label{lesson:MR-C155-bandhne}



\begin{quote}
Before the new one: say the Marathi for to break, then the Marathi for to join.
\end{quote}

\subsection*{You'll want to know: \mr{बांधणे}}
\textbf{\mr{बांधणे}} — \emph{bāndhṇe} — \textquotedblleft{}to tie, to build\textquotedblright{}.

\begin{grammarlens}[title={What you've built}]
One more word for this chapter.
\end{grammarlens}

\subsection*{Your turn}
\begin{itemize}
\raggedright
  \item \emph{Say it:} \emph{bāndhṇe}
  \item \emph{Say it:} \emph{bāndhṇe}, once more
  \item \emph{Say it:} say \emph{bāndhṇe}
  \item \emph{Recall:} say the Marathi for to break, then the Marathi for to join, then say \emph{bāndhṇe} again
\end{itemize}

\begin{tcolorbox}[breakable,colback=teal!4,colframe=teal!35!black,title={Before you move on}]
What does \mr{बांधणे} mean? (\textquotedblleft{}To tie, to build\textquotedblright{}.) Say it once more.
\end{tcolorbox}
\section[śivṇe]{\mr{शिवणे} (śivṇe) — to sew}

\label{lesson:MR-C155-sivne}



\begin{quote}
Before the new one: say the Marathi for to join, then the Marathi for to tie.
\end{quote}

\subsection*{You'll want to know: \mr{शिवणे}}
\textbf{\mr{शिवणे}} — \emph{śivṇe} — \textquotedblleft{}to sew\textquotedblright{}.

\begin{grammarlens}[title={What you've built}]
One more word for this chapter.
\end{grammarlens}

\subsection*{Your turn}
\begin{itemize}
\raggedright
  \item \emph{Say it:} \emph{śivṇe}
  \item \emph{Say it:} \emph{śivṇe}, once more
  \item \emph{Say it:} say \emph{śivṇe}
  \item \emph{Recall:} say the Marathi for to join, then the Marathi for to tie, then say \emph{śivṇe} again
\end{itemize}

\begin{tcolorbox}[breakable,colback=teal!4,colframe=teal!35!black,title={Before you move on}]
What does \mr{शिवणे} mean? (\textquotedblleft{}To sew\textquotedblright{}.) Say it once more.
\end{tcolorbox}
\section[raṅgavṇe]{\mr{रंगवणे} (raṅgavṇe) — to paint}

\label{lesson:MR-C155-rangavne}



\begin{quote}
Before the new one: say the Marathi for to tie, then the Marathi for to sew.
\end{quote}

\subsection*{You'll want to know: \mr{रंगवणे}}
\textbf{\mr{रंगवणे}} — \emph{raṅgavṇe} — \textquotedblleft{}to paint\textquotedblright{}.

\begin{grammarlens}[title={What you've built}]
One more word for this chapter.
\end{grammarlens}

\subsection*{Your turn}
\begin{itemize}
\raggedright
  \item \emph{Say it:} \emph{raṅgavṇe}
  \item \emph{Say it:} \emph{raṅgavṇe}, once more
  \item \emph{Say it:} say \emph{raṅgavṇe}
  \item \emph{Recall:} say the Marathi for to tie, then the Marathi for to sew, then say \emph{raṅgavṇe} again
\end{itemize}

\begin{tcolorbox}[breakable,colback=teal!4,colframe=teal!35!black,title={Before you move on}]
What does \mr{रंगवणे} mean? (\textquotedblleft{}To paint\textquotedblright{}.) Say it once more.
\end{tcolorbox}
\section[kāḍhṇe]{\mr{काढणे} (kāḍhṇe) — to take out}

\label{lesson:MR-C155-kadhne}



\begin{quote}
Before the new one: say the Marathi for to sew, then the Marathi for to paint.
\end{quote}

\subsection*{You'll want to know: \mr{काढणे}}
\textbf{\mr{काढणे}} — \emph{kāḍhṇe} — \textquotedblleft{}to take out\textquotedblright{}.

\begin{grammarlens}[title={What you've built}]
One more word for this chapter.
\end{grammarlens}

\subsection*{Your turn}
\begin{itemize}
\raggedright
  \item \emph{Say it:} \emph{kāḍhṇe}
  \item \emph{Say it:} \emph{kāḍhṇe}, once more
  \item \emph{Say it:} say \emph{kāḍhṇe}
  \item \emph{Recall:} say the Marathi for to sew, then the Marathi for to paint, then say \emph{kāḍhṇe} again
\end{itemize}

\begin{tcolorbox}[breakable,colback=teal!4,colframe=teal!35!black,title={Before you move on}]
What does \mr{काढणे} mean? (\textquotedblleft{}To take out\textquotedblright{}.) Say it once more.
\end{tcolorbox}
